\documentclass[10pt]{article}

\usepackage[margin=0.75in]{geometry}
\usepackage{amsmath,amsthm,amssymb,bm}
\usepackage{xcolor}
\usepackage{cancel}
\usepackage{graphicx}
\usepackage{changepage}
\usepackage{circuitikz}
\usepackage{pgfplots}
\usepackage{minted}
\usepackage{physics}
\usepackage{tikz}
\usepackage{tikz-3dplot}
\usepackage{hyperref}
\usepackage{siunitx}
\usepackage{fontspec}
\usepackage{relsize}
\usepackage{subfig}
\usepackage{todonotes}
\usepackage{contour}
\usepackage{multicol, multirow, booktabs}
\usepackage[breakable]{tcolorbox}
\usepackage[inline]{enumitem}

\theoremstyle{definition}
\newtheorem{problem}{Problem}
\newtheorem{soln}{Solution}

\pgfplotsset{compat=newest}
\usetikzlibrary{lindenmayersystems}
\usetikzlibrary{arrows}
\usetikzlibrary{calc}
\usetikzlibrary{positioning, fit}
\usetikzlibrary{3d, perspective}
\usetikzlibrary{math}
\usetikzlibrary{fadings}
\usetikzlibrary{angles,quotes} 
\usetikzlibrary{decorations.markings}
\usetikzlibrary{tikzmark}

\definecolor{incolor}{HTML}{303F9F}
\definecolor{outcolor}{HTML}{D84315}
\definecolor{cellborder}{HTML}{CFCFCF}
\definecolor{cellbackground}{HTML}{F7F7F7}
\newcommand{\ui}{\hat{i}}
\newcommand{\uj}{\hat{j}}
\newcommand{\uk}{\hat{k}}
\newcommand{\ux}{\hat{\mathbf{x}}}
\newcommand{\uy}{\hat{\mathbf{y}}}
\newcommand{\uz}{\hat{\mathbf{z}}}
\newcommand{\uv}[1]{\hat{\bv{{#1}}}}
\newcommand{\pr}[1]{{#1^\prime}}
\newcommand{\pd}[2][]{{\frac{\partial^{#1}}{\partial {{#2}^{#1}}}}}
\newcommand{\dif}[2][]{{\frac{d^{#1}}{d{{#2}^{#1}}}}}
\newcommand{\grd}{\nabla}

\pgfdeclarelayer{background}  
\pgfsetlayers{background,main}
\AtBeginDocument{\RenewCommandCopy\qty\SI}

\makeatletter
\newcommand{\boxspacing}{\kern\kvtcb@left@rule\kern\kvtcb@boxsep}
\makeatother
\newcommand{\prompt}[4]{
    \ttfamily\llap{{\color{#2}[#3]:\hspace{3pt}#4}}\vspace{-\baselineskip}
}

\newcommand{\thevenin}[2]{
  \begin{center}
    \begin{circuitikz} \draw
      (0,0) -- (2,0) to[battery1, l_=$V_{Th}\eq#1$] (2,2) 
      to[resistor, l_=$R_{Th}\eq#2$] (0,2)
      ;
      \draw [o-] (-.07,2.079);
      \draw [o-] (-.07,0.079);
    \end{circuitikz}
  \end{center}
}

\newcommand{\norton}[2]{
  \begin{center}
    \begin{circuitikz} \draw
      (0,0) -- (3,0) to[american current source, l_=$I_{N}\eq#1$] (3,2) -- (0,2) (2,0)
      to[resistor, l=$R_{N}\eq#2$] (2,2)
      ;
      \draw [o-] (-.07,2.079);
      \draw [o-] (-.07,0.079);
    \end{circuitikz}
  \end{center}
}

\DeclareMathOperator{\Div}{div}
\DeclareMathOperator{\Curl}{curl}
\DeclareMathOperator{\sgn}{sgn}
\DeclareMathOperator{\erfc}{erfc}

\newcommand{\highlight}[1]{\colorbox{yellow}{$\displaystyle #1$}}
\newcommand{\ti}[1]{\widetilde{#1}}
\renewcommand{\div}[1]{\bv{\nabla}\cdot #1}
\renewcommand{\curl}[1]{\bv{\nabla}\times #1}

\newfontface{\Kaufmann}{Kaufmann}
\DeclareTextFontCommand{\kf}{\Kaufmann}
\newcommand{\scriptr}{\fontsize{12pt}{12pt}\kf{r}\fontsize{10pt}{10pt}}

\newfontface{\KaufmannB}{Kaufmann Bold}
\DeclareTextFontCommand{\kfb}{\KaufmannB}
\newcommand{\bscriptr}{\fontsize{12pt}{12pt}\kfb{r}\fontsize{10pt}{10pt}}
\newcommand{\justif}[2]{&{#1}&\text{#2}}
\newcommand{\bv}[1]{\bm{\mathrm{#1}}}

\title{Physics 4700H: Assignment V}
\author{Jeremy Favro (0805980) \\ Trent University, Peterborough, ON, Canada}
\date{\today}

\begin{document}
\maketitle

% PROBLEM 1
\begin{problem}[Kennett problem 2.6] Consider $N$ distinguishable atoms with two energy levels $\pm\epsilon$.
Assume that there are very weak interactions between the atoms such that equilibrium can
be reached but that these can be ignored when calculating equilibrium properties. Let $E$ be
the total energy of the system.
\begin{enumerate}[label=(\alph*)]
  \item What is the entropy of the system as a function of $E$?
  \item What is the temperature of the system as a function of energy?
  \item What are the temperatures corresponding to $U = 0$ and to $U = -N \epsilon$? Comment on
        your results.
\end{enumerate}
\end{problem}
\begin{soln}~
  \begin{enumerate}[label=(\alph*)]
    \item We write $E=N_{+}\epsilon+N_{-}(-\epsilon)=\epsilon\left(N_+-N_-\right)$ where $N_+$ and $N_-$ represent the number of
          particles in the $\epsilon$ and $-\epsilon$ states respectively, and hence $N_++N_-=N$. Such a system has fixed energy (along with volume and number)
          and hence is in the microcanonical ensemble where we may write entropy as $S_B=k_B\ln(\Omega)$ and since there are only two available states,
          $$\Omega=\frac{N!}{\left[\frac{1}{2}\left(N+\frac{E}{\epsilon}\right)\right]!\left[\frac{1}{2}\left(N-\frac{E}{\epsilon}\right)\right]!}$$
          where we've taken the usual $X=N_+-N_-$ and used energy. This gives an entropy of
          $$S_B=k_B\ln(\frac{N!}{\left[\frac{1}{2}\left(N+\frac{E}{\epsilon}\right)\right]!\left[\frac{1}{2}\left(N-\frac{E}{\epsilon}\right)\right]!}).$$
    \item To do this we consider the thermodynamic identity at constant volume,
          $dE=TdS$ which gives us an equation for temperature,
          $$\frac{1}{T}=\left.\frac{\partial S}{\partial E}\right)_{N,V}.$$
          To differentiate the expression we found in (a), unless we want apply the Gamma Function, we must simplify the expression to eliminate the factorials.
          To do this, we break up the $\ln$, then apply Sterling's approximation, $\ln(N!)\approx N\ln(N)-N$,
          \begin{align*}
            \frac{S}{k_B} & =\ln(\frac{N!}{\left[\frac{1}{2}\left(N+\frac{E}{\epsilon}\right)\right]!\left[\frac{1}{2}\left(N-\frac{E}{\epsilon}\right)\right]!})                                           \\
                          & =\ln(N!)-\ln(\left[\frac{1}{2}\left(N+\frac{E}{\epsilon}\right)\right]!)-\ln(\left[\frac{1}{2}\left(N-\frac{E}{\epsilon}\right)\right]!)                                        \\
                          & \approx N\ln(N)-\cancel{N}                                                                                                                                                      \\
                          & \quad-\left[\frac{1}{2}\left(N+\frac{E}{\epsilon}\right)\right]\ln(\frac{1}{2}\left(N+\frac{E}{\epsilon}\right))+\frac{1}{2}\left(\cancel{N}+\cancel{\frac{E}{\epsilon}}\right) \\
                          & \quad-\left[\frac{1}{2}\left(N-\frac{E}{\epsilon}\right)\right]\ln(\frac{1}{2}\left(N-\frac{E}{\epsilon}\right))+\frac{1}{2}\left(\cancel{N}-\cancel{\frac{E}{\epsilon}}\right) \\
                          & =N\ln(N)                                                                                                                                                                        \\
                          & \quad-\left[\frac{1}{2}\left(N+\frac{E}{\epsilon}\right)\right]\ln(\frac{1}{2}\left(N+\frac{E}{\epsilon}\right))                                                                \\
                          & \quad-\left[\frac{1}{2}\left(N-\frac{E}{\epsilon}\right)\right]\ln(\frac{1}{2}\left(N-\frac{E}{\epsilon}\right)).
          \end{align*}
          Now we consider the derivative with respect to $E$,
          \begin{align*}
            \frac{1}{k_B}\frac{\partial S}{\partial E}=\frac{1}{k_BT} & \approx-\frac{\partial }{\partial E}\left[\frac{1}{2}\left(N+\frac{E}{\epsilon}\right)\right]\ln(\frac{1}{2}\left(N+\frac{E}{\epsilon}\right))
            -\frac{\partial }{\partial E}\left[\frac{1}{2}\left(N-\frac{E}{\epsilon}\right)\right]\ln(\frac{1}{2}\left(N-\frac{E}{\epsilon}\right))                                                                    \\
                                                                      & =
            -\frac{1}{2\epsilon}\ln(\frac{1}{2}\left(N+\frac{E}{\epsilon}\right))-\cancel{\left[\frac{1}{2}\left(N+\frac{E}{\epsilon}\right)\right]
            \left(\frac{1}{2\epsilon}\right)\frac{2}{N+\frac{E}{\epsilon}}}                                                                                                                                            \\
                                                                      & \quad +\frac{1}{2\epsilon}\ln(\frac{1}{2}\left(N-\frac{E}{\epsilon}\right))+\cancel{\left[\frac{1}{2}\left(N-\frac{E}{\epsilon}\right)\right]
            \left(\frac{1}{2\epsilon}\right)\frac{2}{N-\frac{E}{\epsilon}}}                                                                                                                                            \\
                                                                      & =\frac{1}{2\epsilon}\left[\ln(\frac{N-\frac{E}{\epsilon}}{N+\frac{E}{\epsilon}})\right]
          \end{align*}
          and so
          $$T(E)=\frac{2\epsilon k_B}{\ln(\frac{N-\frac{E}{\epsilon}}{N+\frac{E}{\epsilon}})}.$$
    \item At $E=0$,
          $$T(E=0)=\frac{2\epsilon k_B}{\ln(\frac{N}{N})}\to \infty$$
          as $\ln(N/N)=\ln(1)=0$.
          At $E=-N\epsilon$,
          $$T(E=-N\epsilon)=\frac{2\epsilon k_B}{\ln(\frac{-2N}{N-N})}\to 0$$
          as $\ln(\infty)\to \infty$. These, somewhat unintuitively due to my desire to make an analogy to a
          kinetic system with a bunch of particles whizzing around, make sense. At $E=-N\epsilon$, the system is in its lowest energy configuration,
          analogous to all the particles not moving, fine. At $E=0$, the system is actually \emph{not} in its lowest energy state, in fact it is in a
          state where nothing is known about the states of the individual particles, a state of maximum entropy.
  \end{enumerate}
\end{soln}

% PROBLEM 2
\begin{problem}
Show that the internal energy may be written in terms of the partition function as
follows,
$$U=k_BT^2\frac{\partial \ln Z}{\partial T}.$$
\end{problem}
\begin{soln}
We know that 
$$Z=\sum_i e^{\epsilon_i/k_B T}.$$
This lets us write that 
$$\frac{\partial \ln Z}{\partial T}=\frac{\partial Z}{\partial T}\frac{\partial \ln (Z)}{\partial Z}=
-\frac{1}{k_B T^2}\sum_i \epsilon_ie^{\epsilon_i/k_B T}
\frac{1}{Z}.
$$
Recalling that $p_i=e^{\epsilon_i/k_B T}/Z$, the expression can be written as
$$\frac{\partial \ln Z}{\partial T}=\frac{1}{k_B T^2}\sum_i \epsilon_ip_i=\frac{U}{k_BT^2}.$$
Multiplying this by $k_BT^2$ we obtain the result,
$$U=k_BT^2\frac{\partial \ln Z}{\partial T}.$$
\end{soln}
\end{document}